\section{Predicting the Span Before Training}
\label{sec:predictfull}

The pilot in Section \ref{sec:predict} asks how far the two-axis statistics generalize. We compute the statistics per dataset and regress them onto the measured best span, evaluated with leave-one-dataset-out cross-validation over fourteen public datasets and two model families (twenty-eight samples: the eight main-suite benchmarks plus the six extension datasets of Appendix \ref{app:ext}). All statistics are computed on the train split only, so no validation or test information enters the features (Appendix \ref{app:stats}). At the dataset level ($n{=}14$) the drift index carries the dominant signal: $\rho = -0.80$ with the log-span ($p = 5\times10^{-4}$) and $\rho = -0.85$ with the long-context benefit ($p = 1\times10^{-4}$). The calendar-aligned autocorrelation is directionally consistent but marginal at this sample size ($\rho = +0.42$ with the log-span, $p = 0.13$; it does not survive Holm correction across the ten statistic$\times$target tests), so we use it as a supporting descriptor rather than an independent axis of evidence. A ridge regressor on the training-free statistics ($k{=}3$ features selected per fold) predicts the held-out log-span with Spearman $\rho = 0.63$ ($p = 4\times10^{-4}$, Figure \ref{fig:predict}) and gets the long-versus-short direction right in 18 of 28 cases. Exact span-class hits (8 of 28) fall below the strongest constant predictor, which always predicts the modal class $L{=}2880$ and hits 14 of 28 because half of the pool is right-censored at the scan ceiling: the regression signal is in the trend and the direction, not the exact tier. The two-feature model ($k{=}2$), which selects exactly the drift and calendar features in eleven of fourteen folds, hits 10 of 28. Long-lag autocorrelation and the detrended spectral score lose significance at this sample size (their earlier signal was carried by the two wide-suite extremes), so we report them only as secondary descriptors. The residual errors concentrate where the two model families prefer different spans on the same data---the part of the phenomenon that dataset-level statistics cannot see; we therefore claim predictability in trend from context statistics, not per series.

\paragraph{Against the standard practice.} The practitioner default is to select the lookback on a validation split. Measured over the 28 dataset--model cells of the pooled suite (fourteen datasets, two model families), validation selection reaches the oracle in 16 of 28 cases with a mean relative regret of 0.92\%, and the statistic predictor in 8 of 28 with 5.33\%---the paired gap favors validation selection (median $+0.9$ percentage points, $t \approx 2.7$, $p \approx 0.01$; Wilcoxon $p \approx 0.02$), so for trained models the zero-cost statistics are a prior on the span, not a substitute for the validation loop. Classical order-selection baselines do not transfer: on the 16-cell main suite (eight datasets, two families) PACF's last-significant-lag saturates at the scan ceiling on all eight datasets, collecting its 6 of 16 exact hits entirely from cells right-censored at the ceiling while its misses are catastrophic (mean relative regret $+$57.8\%); AIC's order is systematically short and never identifies a long-span dataset (2 of 16 exact hits, both on exchange rate). The statistic predictor is the only zero-cost selector that identifies long-span cells (its hits include traffic, the longest-span dataset) and is comparable to AIC on the short side (3 of 16 hits, median regret $+2.3\%$ vs.\ AIC's $+2.9\%$). Where no validation signal exists at all, the same statistics are still decisive: on a frozen foundation model a drift-veto rule matches validation-based span selection on all eight benchmarks tested (Section \ref{sec:method}, Appendix \ref{app:select}).
